\section{Future Research}
\label{sec:future_research}

The first priority for future research is a longer walk-forward evaluation with repeated validation and test windows. More formation periods would permit a stronger assessment of whether the downside-risk ranking, the small ownership-feature gains, and the equal-weight screening result persist across different market conditions. Statistical inference for the portfolio applications would also benefit from a longer return history.

A second extension is to vary the representation of ownership topology while preserving the controlled comparison used in this thesis. Alternative manager-similarity rules, edge definitions, community methods, temporal architectures, and graphs with multiple relation types could capture information omitted by the present construction. Such models should be compared with the same market-only and ownership-based alternatives so that any improvement can be attributed to the added relational information rather than to a broader tuning search.

The prediction task could also be adapted to the limitation identified by the error analysis. Quantile objectives, tail-weighted losses, or classification of drawdowns above a prespecified threshold could focus directly on severe events that the current point forecasts tend to understate. For economic evaluation, future work should use matched investable benchmarks, capacity-aware portfolio construction, and direct estimates of trading, financing, and stock-borrow costs. The near-disappearance of the screening effect under value weighting also warrants explicit analysis of firm size and liquidity before any investment interpretation is advanced.

Overall, the thesis finds that stock-level downside risk can be ranked meaningfully out of sample using information available at the prediction date. Conventional ownership variables add only limited incremental information, and the tested network features and graph models do not deliver a consistent advantage sufficient to justify their complexity. The Fragility Score is therefore best interpreted as a risk-monitoring and portfolio-screening device whose portfolio relevance remains conditional on the investment universe, weighting rule, evaluation period, and implementation assumptions. Stronger claims about trading performance require longer evidence and a more complete treatment of investability and costs.
